\chapter{Methods}\label{chap: Methods}

To simulate the behaviour of fluid motion, it is necessary to integrate the partial differential equations, namely the \acrlong{nse}, numerically with an \fxmarl{a} simulation code.
In this thesis, the high-order \acrfull{sem} \fxmarl{warum ist die abkürzung sem und das wort enthält nur S E?}code Neko is used.
Neko is an open-source \acrfull{cfd} software, that is based on the older Nek5000 \acrshort{cfd} solver originally implemented by~\cite{fischerNek500WebPage}.
The portable framework Neko is written in modern Fortran in an object-oriented manner and is mainly focused on incompressible flow simulations \citep{janssonNekoModernPortable2024}.
It can be run on a \acrfull{cpu} as well as a \acrfull{gpu}.
Only \acrfull{dns} will be performed in this work, however \acrfull{les} models are also implemented in Neko.

\section{Direct Numerical Simulation (DNS)}
\fxmarl{find ich komisch das die abkürzung im titel steht}
In \acrshort{dns}, the \acrlong{nse} are solved without using any turbulence model.
This means, that all scales of turbulent motion are resolved from the largest eddies, down to the viscous scale \fxmarl{s}.
This requires a very fine mesh resolution to perform a \acrshort{dns}, such that the mesh can resolve the smallest eddies.
This leads to increasing resolution requirements in space and time with an increasing Reynolds number.
The requirement for computation time can be estimated as  $\approx Re^3$ for isotropic turbulence and $\approx Re^\frac{37}{14}$ for wall-bounded flows \citep{choiGridpointRequirementsLarge2012}.
The exponential increase in required computation effort makes it obvious that \acrshort{dns} is only viable for low to medium Reynolds numbers and domain sizes.
%This leads to very high computational cost and even with the ever-increasing computational resources of high performance computing, restricts the \acrshort{dns} method to relatively small domain sizes, as well as only low to medium Reynolds numbers.
Nevertheless, \acrshort{dns} is performed regularly in a scientific context, since the accuracy of its results is unmatched by other \acrshort{cfd} simulation methods ZITAT.
In comparison with experimental investigations of turbulent flow \fxmarl{s}, \acrshort{dns} provides significantly more information in the near wall region, since experimental methods often are restricted in the vicinity of the wall.
Especially in rough wall conditions, it can be very challenging or impossible to obtain data about the flow behaviour in the vicinity of roughness features ZITAT.

\subsection{Spectal element method}
The employed \acrshort{sem} method was originally proposed by~\cite{pateraSpectralElementMethod1984} and can be seen as a combination of the \acrfull{fem} and spectral methods.
It achieves significantly higher accuracy and faster convergence than a purely \acrshort{fem} based approach. 
This high-order approach combines the flexibility in spatial discretisation of the \acrshort{fem} with the high accuracy of spectral methods.
The high accuracy of spectral methods is due to the fast spectral convergence.
This is achieved, since the  value of a derivative at a gridpoint is not only dependent by points in its vicinity, but on all points of the domain \fxmarl{komische satzgrammatik/aufbau}.
This leads to spectral convergence, which means that the error decreases faster than exponentially with increasing number of points.
\begin{figure}
    \centering
    \input{plots/Methods/spectral convergence.pgf}
    \caption{Comparison between spectral convergence and algebraic convergence. Solid line representing spectral convergence, dotted line algebraic convergence}
    \label{fig: spectral convergence}
\end{figure}
\autoref{fig: spectral convergence} compares spectral and algebraic convergence. 
Algebraic convergence is typical for other numerical methods, e.g.\ the finite difference method.\fxmarl{warum wird hier nicht FEM genannt sondern noch ein neuer begriff eingebracht?}
The axis scaling is set to log-lin, such that the exponential convergence of spectral methods is representend as a straight line. \fxmarl{finde der satz könnte auch in die caption}

To combine the \acrshort{fem} and spectral methods, first a \acrshort{fem} mesh is constructed. \fxmarl{ich finds iwie etwas komisch dass erst über den fehler gesprochen wird, und dann erst darüber was das eig ist, also aus meienr sicht kann man da einfach die reihenfolge ändern}
\begin{figure}
    \centering
    \input{plots/Methods/SEM_GLL_distribution.pgf}
    \caption{Exemplary \acrshort{sem} element. Subdivided using polynomial order 6.}
    \label{fig: SEM element}
\end{figure}
The mesh is then subdivided into the gridpoints using the \acrfull{gll} distribution.
This places the gridpoints denser at the element boundaries and further apart the element centre.
Additionally, the \acrshort{gll} distribution places points at the element boundaries, which enables continuity across elements.
For each element, the solution is represented by high-order Lagrange interpolating polynomials of degree $N$. \fxmarl{finds blöd das es mit N anfängt und N aufhört}
$N$ is denoting the polynomial order, which means the number of points placed inside each \acrshort{fem} element in each direction.
\fxmarl{ZITAT: im gesamten Absatz}

\subsection{Time discretisation}\label{subsec: time discretization}

When integrating the \acrshort{nse} in time, a feasible time step size must be determined, such that the scheme is stable and produces correct results.
By setting the \acrfull{cfl} number, the necessary time step can be determined.
It's a dimensionless value which is defined with the equation
\begin{equation}\label{eq: cfl number}
    c = \frac{u\Delta t}{\Delta x}
\end{equation}
and is based on the work by~\cite{courantBerPartiellenDifferenzengleiehungen}.
In this formulation $u$ denotes the velocity, $\Delta t$ the current time step and $\Delta x$ the spatial discretisation.
Due to nonuniform meshes and a varying velocity field, the \acrshort{cfl} number is nonuniform over the simulation domain.
If the maximum \acrshort{cfl} number in the domain is above or below the specified value, this formula is rearranged, such that a new time step can be computed.

The time integration of the \acrshort{nse} is done in Neko by an implicit-explicit scheme \citep{janssonNekoModernPortable2024}.
It is based on backward differentiation and extrapolation of the implicit-scheme coefficients.
The viscous diffusion term and time derivative are solved implicitly, while the convective term and body force term are solved explicitly \citep{janssonNekoModernPortable2024}.


\section{Generating rough surfaces}\label{sec:generating rough surfaces}
In real-world engineering applications, surfaces often exhibit irregular roughness.
Previous studies have often used scans of real-world surfaces, to perform simulations, e.g.\ by~\cite{thakkarInvestigationTurbulentFlow2017, demaioDirectNumericalSimulation2023} and~\cite{busseReynoldsnumberDependenceNearwall2017}.
This has the advantage that inherently the real world characteristics of surface roughness is investigated.
This method is, however, restricted by the available surfaces.
Also, the investigated surface parameters cannot be directly controlled; but they are a direct result of the chosen surfaces for scanning.
In the present work, the surface roughness will be generated using an algorithm introduced by~\cite{perez-rafolsGeneratingRandomlyRough2019}, which will be briefly discussed in the following section.
% The following section provides a brief description of the algorithm, followed by an explanation of how the generated surface is implemented in the simulation.

\subsection{Surface generation algorithm}\label{subsec: Surface generation algorithm}

The idea behind the algorithm by~\cite{perez-rafolsGeneratingRandomlyRough2019} is that a realistic surface topography requires both a height distribution and a spatial distribution.
To that extend, two main input parameters need to be specified.
The height distribution of the surface is defined using a probability distribution.
To define the spatial distribution, a power spectrum is used.
The probability distribution and the power spectrum are then used to generate two initial surfaces, $Z_h$ and $Z_s$.
The first surface, $Z_h$, exhibits the desired probability distribution, but not necessarily the power spectrum.
The second surface, $Z_s$, exhibits the power spectrum, but not the probability distribution.
The algorithm then iteratively corrects $Z_h$ and $Z_s$ to match each other, until a specified level of convergence is reached.
The iteration consists of two steps:
\begin{enumerate}
    \item Adjust the $Z_s$ using the formula $\hat{Z}_s^{i+1} = \hat{Z}_h^i \frac{|\hat{Z}_s^i|}{|\hat{Z}_h^i|}$, with the values denoted by $\hat{Z}$ representing their respective Fourier transforms.
    The inverse transform of $\hat{Z}_s^{i+1}$, $Z_s^{i+1}$, shows the specified power spectrum but a changed height probability distribution.
    \item Obtain a $Z_h^{i+1}$ by adjusting the height probability distribution of $Z_s^{i+1}$.
    This is achieved by rank ordering.
\end{enumerate}
In the present work, the Weibull distribution is used as the probability distribution, with the probability density function
\begin{equation}\label{eq: Weibull distribution}
    f(z) = \frac{k}{\lambda} \left(\frac{z}{\lambda}\right)^{k-1} e^{-(z/k)^{k}}
\end{equation}
with $k$ being the shape factor and $\lambda$ representing the scale factor.
The power spectrum used here is defined by
\begin{equation}\label{eq: Power spectrum}
    C(q) = \frac{S_q^2 \beta}{2\pi\left(\beta^2 + q_x^2+q_y^2\right)^{\frac{3}{2}}}
\end{equation}
with $S_q$ being the rms height and $\frac{1}{\beta}$ is the autocorrelation length, as described by~\cite{perez-rafolsGeneratingRandomlyRough2019}.
The specific probability distribution and power spectrum are chosen, since they have been thoroughly tested by~\cite{perez-rafolsGeneratingRandomlyRough2019}.
This shows reasonably fast convergence and low absolute errors are achieved.
It is important to note that the two surfaces $Z_h$ and $Z_s$ are not exactly the same, but similar within the bounds of a defined tolerance.
$Z_h$ will have the desired height probability distribution, but the power spectrum will deviate slightly within the capabilities of the algorithm and the defined tolerances.
Vice versa this holds for $Z_s$.
The surface $Z_h$ is chosen for use in the present work.
\begin{figure}
    \centering
    \begin{subfigure}{0.49\textwidth}
        \subcaption{}
        \includegraphics[width=\linewidth]{Methods/rough_surface_flat.pdf}
        \label{subfig: rough surface flat}
    \end{subfigure}
    \hfil
    \begin{subfigure}{0.49\textwidth}
        \subcaption{}
        \includegraphics[width=\linewidth]{Methods/rough_pipe.pdf}
        \label{subfig: rough pipe}
    \end{subfigure}
    \caption{Comparison between generated rough surface and rough surface applied onto pipe. (a) Plane rough surface. (b) Rough surface mapped onto the pipe geometry}
    \label{fig: comparison between flat and pipe}
\end{figure}

The described algorithm produces a rough plane which needs to be mapped onto the pipe geometry.
\autoref{fig: comparison between flat and pipe} shows both the plane rough surface and the corresponding rough pipe.
The pipe geometry is inherently periodic in the azimuthal direction.
Also, periodic boundary conditions are applied for the inflow and outflow, as described in \autoref{sec: Simulation setup}.
In order to fulfil these conditions, a periodic surface must be generated in length and width direction.
The periodicity in both streamwise and azimuthal direction is given for the surface, since it is a product of Fourier and inverse Fourier transforms, which makes it inherently periodic.

In order to obtain a surface with the desired surface parameters, the correct input values for $k, \lambda, \beta$ and $S_q$ have to be found.
This is achieved by using a quasi-Newton optimisation algorithm implemented in the scipy.optimize python library.
The optimisation algorithm systematically alters the parameters within the parameter space to minimise the target function.
The target function is defined to be the sum of the squared deviations of the considered surface parameters from their respective targets.
To eliminate the random noise from the random number generator of the random distribution, the random seed is fixed throughout this process.



